% problem-000086 2 化学反应与结构综合分析
\section*{2. 化学反应与结构综合分析}
\textit{下列为分问。}

\subsection*{2-1}
关于地球的演化，⽬前主要的观点是，原始地球上没有氧⽓，在⽆氧或氧⽓含量很低时，原核⽣物可利⽤⾃然存在的有机质（⽤ $\mathrm { C H _ { 2 } O }$ 表⽰）和某些⽆机物反应获得能量。例如， $\mathrm { S O } _ { 4 }$ 2−和 $\mathrm { M n O _ { 2 } }$ _{2}可分别作为上述过程的氧化剂，前者⽣成⻩⾊浑浊物，⼆者均放出⽆⾊⽆味⽓体。

\subsection*{2-1}
-1 写出 $\mathrm { S O } _ { 4 } { } ^ { 2 - }$ −参与的释能反应的离⼦⽅程式。

\subsection*{2-1}
-2 写出 $\mathrm { M n O _ { 2 } }$ 参与的释能反应的离⼦⽅程式。

\subsection*{2-2}
约27亿年前光合作⽤开始发⽣，随着这⼀过程的进⾏，地球上氧⽓不断积累，经过漫⻓的演变终达到如今的含量。氧⽓的形成和积累，导致地球上⾃然物质的分布发⽣变化。

\subsection*{2-2}
-1 溶解在海洋中的 $\mathrm { F e ^ { 2 + } }$ 被氧化转变成针铁矿 $\mathrm { \dot { \Omega } a - F e O ( O H ) }$ ⽽沉积，写出离⼦⽅程式。

\subsection*{2-2}
-2 与铁的沉积不同，陆地上的 $\mathrm { M o S } _ { 2 }$ 被氧化⽽导致钼溶于海洋且析出单质硫，写出离⼦⽅程式。

\subsection*{2-3}
2020年8⽉4⽇晚，⻉鲁特港⼝发⽣了举世震惊的⼤爆炸，现场升起红棕⾊蘑菇云，破坏⼒巨⼤。⽬前认定⼤爆炸为该港⼝存储的约2750吨硝酸铵的分解。

另外，亚硝酸铵受热亦发⽣分解甚⾄爆炸，若⼩⼼控制则可⽤在需要充⽓的过程如乒乓球制造之中。

\subsection*{2-3}
-1 写出硝酸铵分解出红棕⾊⽓体的反应⽅程式。

\subsection*{2-3}
-2 写出亚硝酸铵受热分解⽤于充⽓的反应⽅程式。
